\chapter{Là où commence l'aventure}
\label{chap:commence-aventure}

La caravane resterait toujours une maison pour Kael. Elle était l'endroit où il
avait grandi, où il avait appris à lire et à écrire, à tirer à l'arc, à soigner
quelques blessures, à se déplacer sans être entendu et quantité d'autres choses
glanées auprès de ceux qui avaient partagé un morceau de sa route.

Mais Kamatsu avait commencé à suivre sa propre voie et, quelques semaines
après son départ, Kael comprit que le moment était venu de découvrir la sienne.
Au matin, il rassembla ses affaires.

Il passa son arc sur son épaule, prit son paquetage et vérifia presque
machinalement la présence de la dague offerte par Kamatsu à sa ceinture. Puis
il salua ceux qui l'avaient recueilli et vu grandir. Les adieux furent
nombreux. Certains lui souhaitèrent bonne route. D'autres lui donnèrent des
conseils qu'il connaissait déjà ou des recommandations qu'il promit
probablement de suivre. Il y eut quelques accolades, quelques plaisanteries et
davantage d'émotion que Kael n'aurait voulu en montrer.

Puis il regarda une dernière fois les chariots qui avaient constitué son foyer
pendant toutes ces années. La caravane allait continuer sa route.

Kael aussi.

Simplement, ce ne serait plus la même. Ce serait la sienne.

Lorsqu'il quitta les chariots, il ne savait pas exactement où il allait.

Et cela lui convenait parfaitement.

\scenebreak

Durant les trois ou quatre années qui suivirent, Kael voyagea beaucoup sans pour
autant vivre les grandes aventures dont Kamatsu aurait probablement aimé faire
des chansons. Il n'affronta ni dragons ni armées, ne sauva aucun royaume et ne
découvrit aucun trésor oublié.

Sa vie fut beaucoup plus modeste : quelques routes, beaucoup de forêts, des
rencontres de passage et suffisamment de petites mésaventures pour continuer à
apprendre. Kael découvrit surtout qu'il avait besoin de peu tant qu'il restait
dans un environnement qu'il comprenait. Dans les bois, il savait trouver de
l'eau, chasser, reconnaître certaines plantes et trouver un endroit
suffisamment abrité pour passer la nuit.

Lorsqu'une chasse était particulièrement bonne, il gardait ce dont il avait
besoin et vendait le reste à l'étape suivante. Un peu de gibier ou une peau
correctement préparée pouvait devenir quelques pièces, un repas, des
provisions ou quelques flèches. Ailleurs, il échangeait simplement ses
compétences. Il pouvait guider un voyageur, retrouver une piste, participer à
une chasse ou reconnaître un chemin suffisamment sûr pour éviter quelques
mauvaises rencontres.

De temps à autre, il retrouvait même une vie qui lui était particulièrement
familière. Certaines caravanes acceptaient volontiers un archer supplémentaire
capable de surveiller les alentours, de reconnaître les traces laissées près
du camp ou de partir en éclaireur. Kael les accompagnait alors pendant
quelques étapes en échange d'une petite rémunération, du gîte et du couvert.

Le grincement des roues, les feux du soir et les discussions autour des
chariots lui rappelaient immanquablement son enfance. Pendant quelques jours,
parfois quelques semaines, il retrouvait quelque chose qui ressemblait à son
ancienne vie.

Puis venait une intersection, une ville ou simplement un sentier qui attirait
son attention. Kael reprenait son paquetage.

Et continuait seul.

\scenebreak

D'autres travaux étaient plus inhabituels. Dans un village, on lui parla un
jour d'une terrible créature qui terrorisait les environs. À écouter les
habitants, elle était énorme, féroce et probablement responsable d'à peu près
tout ce qui avait disparu au cours des dernières semaines.

Kael accepta de s'en occuper. Il suivit les traces avec beaucoup plus de
prudence que ne le justifia finalement la situation. Le monstre légendaire
s'avéra être une bête sauvage certes inhabituellement grande et devenue
dangereuse, mais très loin de la créature fantastique décrite à la taverne.

Il s'en débarrassa néanmoins et reçut la petite prime promise. Il apprit
surtout qu'entre une créature rencontrée dans une forêt et la même créature
racontée le soir autour de trois chopes, quelques détails avaient parfois
tendance à apparaître.

Ces années lui permirent d'acquérir de l'expérience, mais Kael était encore
jeune. Il apprenait toujours. Il commettait des erreurs, surestimait parfois
ses capacités et connaissait davantage les dangers ordinaires des routes que
les véritables menaces dont parlaient les légendes.

Son histoire d'aventurier, au fond, n'avait pas encore réellement commencé.

\scenebreak

Kamatsu n'avait pas disparu de sa vie pour autant. Leurs chemins étaient
désormais différents, mais il arrivait qu'ils se croisent. Deux ou trois fois
au cours de ces années, Kael retrouva son frère dans une taverne, parfois
presque par hasard, parfois parce que l'un avait appris que l'autre se trouvait
dans les environs.

Kamatsu voyageait lui aussi, allant là où ses histoires, ses chansons ou
simplement sa curiosité le conduisaient. Ces retrouvailles pouvaient durer une
soirée ou quelques jours.

Autour d'une table, chacun racontait ce qu'il avait vécu depuis leur dernière
rencontre. Kamatsu avait de nouvelles histoires ; Kael de nouvelles routes,
quelques chasses et généralement une ou deux anecdotes qu'il considérait
parfaitement ordinaires jusqu'à ce que son frère décide qu'elles méritaient
manifestement d'être racontées autrement.

Kael découvrit ainsi que certaines de ses aventures devenaient beaucoup plus
impressionnantes après être passées entre les mains de Kamatsu.

Il protesta parfois.

Cela n'eut généralement aucun effet.

Puis ils se séparaient de nouveau. Ils savaient désormais que cela ne
signifiait jamais vraiment se quitter.

\scenebreak

Quelques années après avoir quitté la caravane, les routes conduisirent
finalement Kael jusqu'à Waterdeep.

Il y découvrit rapidement une faiblesse assez importante dans son mode de vie :
il était beaucoup plus facile d'être pauvre dans une forêt que dans une ville.
Dans les bois, ne pas posséder une seule pièce n'était pas particulièrement
inquiétant. L'eau ne demandait pas à être payée, dormir sous un arbre non plus
et, avec suffisamment de patience, il était généralement possible de trouver
quelque chose à manger.

À Waterdeep, presque tout semblait avoir un propriétaire.

Et, surtout, un prix.

Kael aurait évidemment pu chasser. Il comprit néanmoins assez rapidement que
chasser les animaux domestiques des habitants n'était probablement pas la
meilleure manière de financer son séjour. Il choisit donc une solution qui lui
paraissait nettement plus raisonnable : vendre ses services.

Kael aborda quelques personnes et proposa ce qu'il savait faire. Il pouvait
servir de guide hors de la ville, accompagner un voyageur sur une route
dangereuse, retrouver une piste, participer à une chasse ou protéger un petit
convoi. Il ne réclamait même pas nécessairement beaucoup d'argent. Quelques
pièces ou un repas contre un travail honnête lui auraient parfaitement
convenu.

Le problème fut que les gardes de Waterdeep n'interprétèrent pas exactement la
situation de cette manière. De leur point de vue, un jeune Orc sans argent qui
abordait les passants dans les rues en espérant obtenir quelques pièces
ressemblait fortement à un mendiant.

Kael tenta d'expliquer qu'il ne demandait justement pas qu'on lui donne de
l'argent. Il proposait de travailler en échange.

La différence lui paraissait importante.

Aux gardes, beaucoup moins.

La nouvelle Lord de Waterdeep, Jeanne Von Brocht, ne voulait pas de mendiants
dans sa ville. Et puisque Kael semblait correspondre suffisamment à cette
définition aux yeux de ceux chargés d'appliquer ses ordres, ses explications ne
changèrent pas grand-chose à son sort.

Il fut arrêté.

\scenebreak

Quelques heures plus tard, assis derrière les barreaux, Kael eut tout le loisir
de réfléchir à l'absurdité de la situation. Il avait vécu seul dans les forêts,
affronté des créatures sauvages, traversé des régions inconnues et accompagné
des caravanes sur des routes dangereuses.

Et pourtant, c'était en essayant honnêtement de gagner de quoi acheter son
prochain repas qu'il avait réussi à finir en prison.

Les villes étaient décidément des endroits beaucoup plus dangereux qu'elles
n'en avaient l'air.

Kael ignorait encore que cette mésaventure mettait un terme à ses quelques
années d'errance solitaire.

Car il n'était pas seul dans cette prison.

D'autres personnes avaient été jetées derrière les mêmes barreaux, chacune avec
son histoire, ses raisons d'être là et probablement sa propre opinion sur la
nouvelle Lord de Waterdeep. Kael ne connaissait encore aucune d'entre elles.

Mais cela n'allait pas durer.

Il était parti quelques années auparavant pour trouver sa voie. Jusqu'ici, elle
l'avait conduit à travers des forêts, des villages, quelques tavernes,
plusieurs caravanes et, finalement, d'une manière assez peu glorieuse, jusqu'au
fond d'une cellule.

\begin{center}
\textbf{C'est pourtant derrière ces barreaux que son véritable voyage allait
commencer.}
\end{center}